\documentclass[10pt,a4paper]{amsart}
\usepackage[margin=19mm]{geometry}
\usepackage{amsmath,amssymb,mathtools}
\usepackage[scr=rsfs,cal=euler]{mathalfa}
\usepackage{xfrac}
\usepackage[unicode,hidelinks,bookmarks=false]{hyperref}
\newcommand{\ie}{i.e.\ }
\newcommand{\cf}{cf.\ }
\newcommand{\resp}{respectively\ }
\newcommand{\Subsection}{\subsection}
\providecommand{\nobreakdash}{\nobreak\hbox{-}}
\setlength{\emergencystretch}{2em}
\multlinegap0pt
\usepackage{siunitx}
\usepackage{siunitx}
\usepackage{siunitx}
\usepackage{siunitx}
\usepackage{amssymb}
\usepackage{mathtools}
\usepackage{tikz}
\usepackage{xcolor}
\usepackage{siunitx}
\usepackage[shortlabels]{enumitem}
\newtheorem{proposition}{Proposition}
\newcommand\bC{{\mathbb C}}
\newcommand{\sH}{{\mathcal H}}
\title{Numerical Elimination: Computing Complements\\ of Real Hypersurfaces Using Pseudo-Witness Sets}
\author{Paul Breiding, John Cobb, Aviva K. Englander, Nayda Farnsworth, Jonathan D. Hauenstein, Oskar Henriksson, David K. Johnson, Jordy Lopez Garcia, Deepak Mundayur}
\date{Source: arXiv:2601.04383v3 (CC BY 4.0)}
\begin{document}
\maketitle

\noindent\emph{Source and licence.} Numerical Elimination: Computing Complements\\ of Real Hypersurfaces Using Pseudo-Witness Sets. Version arXiv:2601.04383v3 (CC BY 4.0), distributed by arXiv under the Creative Commons Attribution 4.0 International licence, http://creativecommons.org/licenses/by/4.0/, as recorded in the arXiv licence metadata for this exact version and shown on the abstract page https://arxiv.org/abs/2601.04383v3. Full licence notice: \texttt{licenses/arxiv-2601.04383-CC-BY-4.0.txt}.\par\smallskip

\noindent\textit{Editorial scope.} Counted unit: the article's proposition prop:gradient (source0.tex lines 703--710). The article's complete original proof of it is delivered with it (source0.tex lines 712--722, one \texttt{proof} environment of 11 lines). No mathematics is written, repaired or re-derived: the statement and the proof are the article's own lines, in the article's own notation, and no retained line is substituted or re-flowed.  Retained uncounted as the source-local context the proof needs: lines 698--700.  Nothing was omitted from any retained span, so the package carries no omission record.  No retained line contains a citation, so the wrapper carries no bibliography and no bibliography entry.  No retained line contains a figure environment, an image-inclusion command or drawing code, so no image is delivered and the package declares no asset dependency.  Editorial formatting only, in addition: an AMS \texttt{amsart} preamble that restates the article's own declarations and macros used by the retained lines, namely the environments \texttt{proposition} on separate counters, together with the standard packages those lines need.  The retained environments print in this document as Proposition 1 in order of appearance, whereas the article prints the same items with the numbering of its own full document.  The complete native source of the article as distributed in the arXiv e-print for this exact version is bundled unchanged as \texttt{sources/192291/source0.tex}.
\begin{equation}\label{eq:directional-rational-h-complex}
    \frac{\frac{\mathrm{d}}{\mathrm{d}t}h(p+tb)}{h(p+tb)} = \sum_{j=1}^{\deg h}\frac{1}{t-t_{j}},
\end{equation}

\begin{proposition}
\label{prop:gradient}
For $p\in \bC^{k}\smallsetminus\sH$ and a generic direction $b \in \bC^{k}$,
\begin{equation*}
    \frac{\nabla h(p)}{h(p)} \cdot b = -\sum_{j=1}^{\deg h}s_{j}(p,b).
\end{equation*}
When $p$ and $b$ are real, the left-hand side agrees with the directional derivative of $\log \lvert h\rvert$ at $p$ in direction $b$.
\end{proposition}

\begin{proof}
Since 
\begin{equation*}
\frac{\mathrm{d}}{\mathrm{d}t} h(p +tb)\vert_{t=0} = \nabla h(p)\cdot b,
\end{equation*}
evaluating~\eqref{eq:directional-rational-h-complex} at $t=0$ gives
\begin{equation*}
\frac{\nabla h(p)}{h(p)}\cdot b = -\sum_{j=1}^{\deg h}\frac{1}{t_{j}} = -\sum_{j=1}^{\deg h}s_{j}(p,b),
\end{equation*}
where the last equality follows from $s_j = 1/t_j$. When $p$ and $b$ are real, the last statement follows from $\nabla \log\lvert h\rvert = (\nabla h)/h$.
\end{proof}

\end{document}
